\bigskip
\bigskip
\subsubsection*{Խնդիրներ և հարցեր թեմա 13-ի վերաբերյալ}

\begin{enumerate}[label=\thesection.\arabic*.]

% 13.1
\item Ճի՞շտ է արդյոք, որ կամայական $X$ և $Y$ տարածությունների դեպքում \linebreak $P_X: X \times Y \to X$ կանոնական պրոյեկցիան փակ արտապատկերում է:

\begin{hint}
Տե՛ս \hyperref[ch9ex4]{օրինակ 4}-ը թեմա 9-ում:
\end{hint}

% 13.2
\item Ապացուցեք․ կամայական $X$ և $Y$ տարածությունների դեպքում $X \times Y$-ը հո\-մե\-ո\-մորֆ է $Y \times X$-ին:

\begin{hint}
Հիմնավորեք, որ $f: X \times Y \to Y \times X$, $f(x,y)=(y,x)$ արտա\-պատ\-կերումը հոմեոմորֆիզմ է:
\end{hint}

% 13.3
\item Ապացուցեք․ կամայական $X, Y, Z$ տարածությունների դեպքում $(X \times Y) \times Z$-ը հոմեոմորֆ է $X \times (Y \times Z)$-ին:

\begin{hint}
Հիմնավորեք, որ $g: (X \times Y) \times Z \to X \times (Y \times Z)$, $g\big((x,y),z\big)=\big(x,(y,z)\big)$ արտապատկերումը հոմեոմորֆիզմ է:
\end{hint}

% 13.4
\item Գոյություն ունի՞ $(X, \tau)$ տարածություն, որը հոմեոմորֆ է իր $(X \times X,\, \tau \times \tau)$ քառակուսուն:

\begin{hint}
Հիմնավորելով, որ $X$-ը պետք է լինի անվերջ բազմություն՝ դիտ\-ար\-կեք, օրինակ, բնական թվերի բազմությունը դիսկրետ տոպոլոգիայով:
\end{hint}

% 13.5
\item Ճի՞շտ է արդյոք, որ կամայական $X$ անվերջ բազմություն դիսկրետ տոպոլո\-գի\-ա\-յով հոմեոմորֆ է իր $(X \times X,\, \text{դիսկր.} \times \text{դիսկր.})$ քառակուսուն:

% 13.6
\item Դիտարկենք $A=[-1,1] \times [-1,1]$ քառակուսին $\R^2$ հարթությունում և նրան ներգծված $B=\{(x_1, x_2);\ x_1^2+x_2^2 \le 1\}$ շրջանը: Ապացուցեք, որ դրանք հոմեո\-մորֆ տարածություններ են:

\begin{hint}
Նախ համոզվեք, որ $A$-ի եզրագծի ամեն մի $(x_1, x_2)$ կետի դեպքում կամ $|x_1|=1, |x_2| \le 1$, կամ $|x_2|=1, |x_1| \le 1$։ Հիմնվելով դրա վրա՝ ապացուցեք․
\[
h: A \to B, \qquad h(x_1, x_2) = \frac{\max\{|x_1|, |x_2|\}}{\sqrt{x_1^2+x_2^2}}(x_1, x_2)
\] 
արտապատկերումը հոմեոմորֆիզմ է:
\end{hint}

% 13.7
\item Ապացուցեք․ $\R^n$ էվկլիդեսյան մետրիկային տարածության $\R^n \setminus \{0\}$ ենթա\-տա\-րածությունը հոմեոմորֆ է $S^{n-1} \times \R$ արտադրյալին, որտեղ $0$-ն $\R^n$-ի սկզբնա\-կետն է, իսկ $S^{n-1}$-ը $\big\{x=(x_1, \dots, x_n);\ \|x\|^2 = \sum\limits_{i=1}^n x_i^2=1\big\}$ սֆերան է:

\begin{hint}
Հիմնավորեք, որ 
\[
f: \R^n \setminus \{0\} \to S^{n-1} \times \R, \qquad f(x)=\left(\frac{1}{\|x\|}(x_1, \dots, x_n),\ \lg \|x\|\right)\]
արտապատկերումը հոմեոմորֆիզմ է:
\end{hint}

% 13.8
\item Կամայական $f: X \to Y$ արտապատկերման համար սահմանվում է նրա գրա\-ֆիկը որպես $X \times Y$ արտադրյալի $\Gamma_f = \{(x,y);\, y=f(x)\}$ ենթաբազ\-մու\-թյուն։ Ապացուցեք․
\begin{enumerate}
    \item[ա)] եթե $f$-ը անընդհատ է, ապա $\Gamma_f$-ը որպես $X \times Y$-ի ենթատարածություն հոմեոմորֆ է $X$-ին;
    \item[բ)] եթե $X$-ը հաուսդորֆյան տարածություն է, ապա $\Gamma_f$-ը փակ ենթաբազմու\-թյուն է $X \times Y$-ում:
\end{enumerate}

% 13.9
\item Ապացուցեք. $X$-ը հաուսդորֆյան տարածություն է այն և միայն այն դեպքում, երբ $X \times X$-ի $\Delta = \{(x,x);\, x \in X\}$ անկյունագիծը փակ ենթաբազմություն է $X \times X$-ում:

\begin{hint}
Օգտվեք նախորդ խնդրի բ) պնդումից՝ որպես $f$ վերցնելով $X$-ի նույնական արտապատկերումը:
\end{hint}

\end{enumerate}

\par Մինչև հաջորդ խնդիրներին անցնելը տանք մի սահմանում։ Ասում են, որ տո\-պո\-լոգիական տարածություններին առնչվող տվյալ \textbf{հատկությունը պահպանվում է ուղիղ արտադրյալի դեպքում}, եթե ամեն անգամ, երբ որևէ երկու տարածություն\-ներ բավարարում են այդ հատկությանը, նրանց ուղիղ արտադրյալը նույնպես բա\-վա\-րարում է այդ հատկությանը։
\par Օրինակ՝ \hyperref[ch13thm7]{թեորեմ 7}-ից հետևում է, որ անջատելիության $\textrm{T}_2$ աքսիոմը պահպան\-վում է ուղիղ արտադրյալի դեպքում։ Մյուս կողմից որպես տեղեկություն նշենք, որ նոր\-մա\-լու\-թյունը (որպես անջատելիության աքսիոմ) չի պահպանվում ուղիղ արտա\-դրյալի դեպքում:

\begin{enumerate}[label=\thesection.\arabic*., resume]

% 13.10
\item Ապացուցեք, որ $\textrm{T}_0$, $\textrm{T}_1$ աքսիոմները պահպանվում են ուղիղ արտադրյալի դեպ\-քում:

% 13.11
\item Ապացուցեք, որ սեպարաբելությունը պահպանվում է ուղիղ արտադրյալի դեպ\-քում:

% 13.12
\item Դիցուք $B_1=\{U_i;\, i \in I\}$ և $B_2=\{V_j;\, j \in J\}$ ընտանիքները տոպոլոգիայի բազա են համապատասխանաբար $(X, \tau)$ և $(Y, \sigma)$ տարածությունների հա\-մար։ Ապացուցեք․ $B=\{U_i \times V_j;\, i \in I, j \in J\}$ ընտանիքը $(X \times Y, \tau \times \sigma)$ տարածու\-թյան տոպոլոգիայի բազա է։

% 13.13
\item Ապացուցեք․ հաշվելիության երկրորդ աքսիոմը պահպանվում է ուղիղ ար\-տա\-դրյալի դեպքում:

% 13.14
\item Ապացուցեք․ հաշվելիության առաջին աքսիոմը պահպանվում է ուղիղ ար\-տա\-դրյալի դեպքում:

\begin{hint}
Նախ ձևակերպեք և ապացուցեք 13.12 խնդրին նմանակ պնդում կետի շրջակայքերի բազա տերմինով:
\end{hint}

\end{enumerate}
